% Chapter 21: The Erd\H{o}s Discrepancy Problem
\let\LocalMacro\relax

\chapter{The Erdős Discrepancy Problem}

\section{The Problem}

\subsection{Informal}
Imagine assigning either plus one or minus one to every counting number: 1, 2, 3, 4, and so on forever. Now look at every evenly-spaced pattern you can find---the even numbers, the multiples of 3, the multiples of 7, and so on---and add up the signs along each pattern. Can you assign the signs so that \emph{every} such sum stays small, no matter how far out you go?

\begin{historicalbox}
Paul Erdős posed this question in 1932 and offered \$500 for a solution. The problem resisted for 83 years---it is one of the simplest-to-state, hardest-to-solve puzzles in all of mathematics. Erdős called it one of his favorites.
\end{historicalbox}

\subsection{Formal}


\section{Solution}

Who solved it and what techniques were required.

\section{Mathematical Theory}


\subsection{Prerequisites}
This chapter assumes familiarity with:
\begin{itemize}
\item Basic analysis (sequences, series, $\limsup$)
\item Elementary combinatorics
\item Fourier analysis on finite abelian groups
\item Basic probability theory
\end{itemize}

\subsection{Discrepancy and Multiplicative Functions}
A sequence $x: \mathbb{N} \to \{-1, +1\}$ is \emph{completely multiplicative} if $x(mn) = x(m)x(n)$. For such sequences, the discrepancy simplifies:
\[
\sum_{k=1}^N x_{kd} = x_d \sum_{k=1}^N x_k
\]
so the discrepancy is controlled by the partial sums of the original sequence.

\begin{conjecture}[Erdős Discrepancy Conjecture]
For any sequence $x_n \in \{-1, +1\}$, the discrepancy $D = \infty$.
\end{conjecture}

\subsection{Connection to the Liouville Function}
The Liouville function $\lambda(n) = (-1)^{\Omega(n)}$ (where $\Omega(n)$ counts prime factors with multiplicity) is completely multiplicative. The Erdős discrepancy conjecture for $\lambda$ is equivalent to the \emph{prime number theorem in arithmetic progressions} (which is known), but the general conjecture is much stronger.

\subsection{Previous Approaches}
\begin{itemize}
\item \textbf{Beck-Fiala theorem \cite{beckfiala1981}}: For set systems with bounded degree, discrepancy $O(\sqrt{d} \log n)$.
\item \textbf{Spencer's theorem \cite{spencer1985}}: Six standard deviations suffice ($O(\sqrt{n})$).
\item \textbf{Banaszczyk's theorem (1998)}: Discrepancy $O(\sqrt{\log n})$ for arbitrary vectors.
\item \textbf{Computer searches}: Long sequences with low discrepancy were found (length 1124 with discrepancy 2), suggesting the conjecture might be false.
\end{itemize}

\section{The Discovery}

\subsection{Tao's Breakthrough \cite{tao2016discrepancy}}
Terence Tao announced a proof on September 18, 2015, using a combination of:

\begin{enumerate}
\item \textbf{Entropy decrement argument}: A new averaging technique showing that if discrepancy is bounded, the sequence must correlate with a Dirichlet character.
\item \textbf{Elliott's conjecture}: A generalization of Chowla's conjecture on correlations of multiplicative functions, proved by Tao using entropy methods.
\item \textbf{Fourier analysis and nilsequences}: Showing that bounded discrepancy forces correlation with a nilsequence, which then forces unbounded discrepancy via Gowers norms.
\end{enumerate}

\begin{theorem}[Tao \cite{tao2016discrepancy}]
For any sequence $x_n \in \{-1, +1\}$, the discrepancy is infinite.
\end{theorem}

\begin{proof}[Sketch]
Assume for contradiction that $D < \infty$. Then:
\begin{enumerate}
\item By the entropy decrement argument, $x$ correlates with a Dirichlet character $\chi$ of bounded conductor.
\item By Elliott's conjecture (proved via entropy), this correlation forces $x$ to behave like $\chi$.
\item But a completely multiplicative function of bounded discrepancy must be close to a Dirichlet character, and such characters have unbounded discrepancy along some arithmetic progression (by Dirichlet's theorem on primes in AP).
\item Contradiction.
\end{enumerate}
\end{proof}

\begin{highlightbox}
Tao's proof was remarkable because it combined probabilistic combinatorics (entropy), analytic number theory (Elliott's conjecture), and ergodic theory (nilsequences). The proof is non-constructive and does not give explicit bounds on how fast the discrepancy grows.
\end{highlightbox}

\subsection{Subsequent Developments}
\begin{itemize}
\item \textbf{Polymath project}: Verified and streamlined the proof.
\item \textbf{Quantitative bounds}: No explicit growth rate for discrepancy is known (the proof gives a tower-of-exponentials type bound).
\item \textbf{Variants}: Erdős discrepancy for complex units, for sequences in $\{-1,0,1\}$, etc.
\end{itemize}

\section{Impact}
\begin{itemize}
\item \textbf{Combinatorics}: First major application of entropy methods to discrepancy theory.
\item \textbf{Number theory}: Elliott's conjecture (proved en route) is a major result on correlations of multiplicative functions.
\item \textbf{Computer science}: Implications for online algorithms and streaming algorithms.
\item \textbf{Erdős legacy}: Erdős's \$500 prize was finally claimed (Tao donated it to charity).
\end{itemize}

\section{Exercises}

\begin{exercise}[Basic]
Show that the sequence $x_n = 1$ for all $n$ has discrepancy $D = \infty$.
\end{exercise}

\begin{exercise}[Basic]
Show that any periodic sequence $x_n \in \{-1, +1\}$ has infinite discrepancy.
\end{exercise}

\begin{exercise}[Intermediate]
Let $x_n$ be completely multiplicative with $x_p = 1$ for all primes $p \equiv 1 \pmod{4}$ and $x_p = -1$ for $p \equiv 3 \pmod{4}$. Show that $x_n$ has infinite discrepancy.
\end{exercise}

\begin{exercise}[Challenge]
Explain the entropy decrement argument in your own words: how does averaging over large primes reduce the entropy of the sequence?
\end{exercise}

\section{Timeline}

\begin{tabularx}{\linewidth}{lX}
\toprule
\textbf{Year} & \textbf{Event} \\ \midrule
1932 & Erdős poses the discrepancy problem \\ 
1950s & Erdős offers \$500 prize \\ 
1981 & Beck-Fiala theorem \\ 
1985 & Spencer's "six standard deviations" theorem \\ 
2015 & Tao proves the Erdős discrepancy conjecture \\ 
2016 & Polymath project streamlines the proof \\ \bottomrule
\end{tabularx}

\section{Citations}

\begin{enumerate}
\item Tao, T. (2016). ``The Erdős discrepancy problem.'' Discrete Analysis, 1, 29 pp.
\item Tao, T. (2016). ``The logarithmically averaged Chowla and Elliott conjectures for two-point correlations.'' Forum of Mathematics, Pi, 4, e8.
\item Erdős, P. (1957). ``Some unsolved problems.'' Michigan Mathematical Journal, 4(3), 291--300.
\item Beck, J., \& Fiala, T. (1981). ``Integer-making theorems.'' Discrete Applied Mathematics, 3(1), 1--8.
\item Spencer, J. (1985). ``Six standard deviations suffice.'' Transactions of the AMS, 289(2), 679--706.
\end{enumerate}
